\cotefile{79}
\scanpage{160}
\mspage{81}

\uline{IX} \uline{Graphes à \deleted{opérateurs}\inserted{homothéties}.}

Nous \deleted{\unsure{prendr|prendrons}} appelons ici structure de \uline{graphes} sur un ens. de sommets $S$, la donnée d'une partie

\[ A\subset\unsure{\mathcal P|\mathfrak P}_2(S) \tag*{(1)} \]

\deleted{\unsure{N}}Nous appelerons \uline{graphe à \deleted{opérateurs}\inserted{homothéties}} une structure sur un ens. de base $S$ (l'ens de sommets) définie par la donnée de deux structures

\[ A_0\subset\mathcal P_2(S),\qquad \underset{\text{\uline{sous-groupe}}}{Z\subset\unsure{\mathfrak S|\mathscr S}_S} \tag*{(2)} \]

satisfaisant aux conditions qui suivent. Notons tout de suite qu'on ne suppose nullement que \deleted{\unsure{ceci} \uline{GH1} $Z$ est un sous-groupe de $\mathfrak S_S$} les \deleted{$\unsure{g}$} $\lambda\in Z$ respectent la structure de graphe \deleted{$\unsure{A_\lambda}$} $A$ de $S$. On suppose par contre :

\[
\begin{aligned}
 &\forall(s,t)\in S,\ \lambda\in Z,\quad \underline{\exists f}\in\mathfrak S_S,\ \text{centralisant }Z,\ \text{tel que}\\
 &\{s,t\}\in A_0\Longleftrightarrow\{f^{-1}\deleted{\unsure{\lambda|1}}s,f_{\lambda}^{-1}\deleted{\unsure{1|\lambda}}t\}\in A_0
\end{aligned}\tag*{\shortstack{(3) \deleted{\unsure{G}}\\(GH1)}}
\]

Cette condition peut se décomposer en trois

\[
\mathrm{GH1}'\left\{\begin{aligned}
\text{a)}\quad &\forall\lambda\in Z,\ \exists A_\lambda\in\mathcal P(\mathcal P_2(S))\quad\text{structure de graphe}\\
 &\text{sur }S\ \text{(évidemment unique) telle que }\forall s,t\in S\\
 &\{s,t\}\in A_0\Longleftrightarrow\{s,\lambda t\}\in A_\lambda\\
 &\text{\deleted{$[$ NB les \unsure{arêtes|étoiles} de $(S,A_0)$ et de $(S,A_\lambda)$ sont les mêmes $]$}}\\
\text{b)}\quad &\text{Cette structure de graphe $A_\lambda$ est isomorphe à $A_0$.}\\
 &\text{\deleted{\unsure{ou mieux}}}\\
 &\text{\deleted{c) Parmi les \unsure{isom.} $f\in\mathfrak S_S$ qui \unsure{établiss|centralisent} cet is}}\\
\text{c)}\quad &\text{\deleted{\unsure{isomorphisme, il en existe}} \inserted{On veut qu'il existe} un isom$^{(f)}$ de}\\
 &\text{la structure $(S,A_0,Z)$ sur la structure}\\
 &\text{$(S,A_\lambda,Z)$, et plus précisément un isom}\\
 &\text{qui \deleted{induit} \inserted{induise} l'identité sur $Z$, i.e.}\\
 &\text{par un $f\in\mathfrak S_S$ qui centralise $Z$.}
\end{aligned}\right.\tag*{(4)}
\]

% slanted
\margin{\deleted{Ceci}\inserted{a)} implique que si $\lambda,\mu\in Z$, alors $\forall s\in S$, on a $\mathrm{Et}_{A_0}((\lambda\mu)s)=\mathrm{Et}_{A_0}((\mu\lambda)s)$ donc $\lambda\mu=\mu\lambda$ si on suppose que $\begin{array}{rcl}S&\longrightarrow&\mathcal P(S)\\s&\longmapsto&\mathrm{Et}_A(s)\end{array}$ est injective. En tout état de cause, il est raisonnable de supposer $Z$ commutative.} (5) 

\deleted{On trouve} \inserted{On pose}

\[ \pi=\{A_\lambda\mid\lambda\in Z\}\subset\mathcal P(\mathcal P_2(S)) \]

\deleted{et on fait abstraction des \unsure{liens}} \inserted{et pour $i\in\pi$ on note $A_i$ au lieu de $A_\lambda$.}

% archive sheet number: 159
